\documentclass[11pt]{article}

\usepackage[T1]{fontenc}
\usepackage{lmodern}
\usepackage{amsmath,amssymb,amsthm,mathtools}
\usepackage[margin=1in]{geometry}
\usepackage[hidelinks]{hyperref}

\newtheorem{theorem}{Theorem}
\newcommand{\R}{\mathbb{R}}
\newcommand{\C}{\mathbb{C}}
\newcommand{\norm}[1]{\left\lVert #1\right\rVert}

\title{Negative Constrained Energy at Every Positive Mass\\
for a Fourth-Order Cubic--Quintic Hamiltonian on $\R^2$}
\author{Richard J. Reyes}
\date{October 26, 2026}

\begin{document}
\maketitle

\begin{abstract}
We prove that the constrained energy infimum of a fourth-order
focusing-cubic/defocusing-quintic Hamiltonian on $\R^2$ is strictly
negative at every positive mass. An anisotropic Gaussian modulated by
a plane-wave carrier yields an exact two-term expression in the
spatial scale. Choosing sufficient anisotropy makes the leading
coefficient negative. A pointwise square-completion identity also
provides a uniform lower bound. This is a variational infimum result;
it does not by itself prove existence of minimizers or orbital stability.
\end{abstract}

\section{Statement}
Let $\beta,\gamma,k>0$, and consider the functional
\begin{equation}\label{eq:H}
 H(u)=\norm{(\Delta+k^2)u}_{L^2(\R^2)}^2
 -\frac{\beta}{2}\norm{u}_{L^4(\R^2)}^4
 +\frac{\gamma}{3}\norm{u}_{L^6(\R^2)}^6,
 \qquad u\in H^2(\R^2;\C).
\end{equation}
For each $m>0$, define
\begin{equation}\label{eq:I}
 I(m)=\inf\bigl\{H(u):u\in H^2(\R^2;\C),\ \norm{u}_2^2=m\bigr\}.
\end{equation}

\begin{theorem}\label{thm:negativity}
For every $\beta,\gamma,k>0$ and every $m>0$,
\begin{equation}\label{eq:main}
 \boxed{-\frac{3\beta^2}{16\gamma}\,m\ \le\ I(m)\ <\ 0.}
\end{equation}
In particular, the constrained energy infimum is strictly negative at
arbitrarily small positive mass.
\end{theorem}

\section{Proof}
\begin{proof}
\textbf{Step 1: A normalized anisotropic Gaussian.}
For $a,b>0$, set
\begin{equation}\label{eq:f}
 f_{a,b}(x,y)=\frac{1}{\sqrt{\pi ab}}
 \exp\!\left(-\frac{x^2}{2a^2}-\frac{y^2}{2b^2}\right).
\end{equation}
A direct Gaussian integral gives $\norm{f_{a,b}}_2^2=1$.
For $R>0$, introduce the Gaussian carrier trial
\begin{equation}\label{eq:trial}
 u_R(x,y)=\frac{\sqrt{m}}{R}
 f_{a,b}\!\left(\frac{x}{R},\frac{y}{R}\right)e^{ikx}.
\end{equation}
It belongs to $H^2(\R^2;\C)$, and the change of variables
$(X,Y)=(x/R,y/R)$ shows that
\begin{equation}\label{eq:mass}
 \norm{u_R}_2^2
 =m\int_{\R^2}|f_{a,b}(X,Y)|^2\,dX\,dY=m.
\end{equation}
Thus every $u_R$ is admissible in \eqref{eq:I}.

\medskip
\noindent\textbf{Step 2: The fourth-order kinetic term.}
Write $u_R=e^{ikx}v_R$, where the real-valued envelope is
\[
 v_R(x,y)=\frac{\sqrt m}{R}f_{a,b}\left(\frac{x}{R},\frac{y}{R}\right).
\]
Then
\begin{equation}\label{eq:carrier}
 (\Delta+k^2)u_R
 =e^{ikx}(\Delta v_R+2ik\,\partial_x v_R).
\end{equation}
The two summands inside the parentheses are respectively real and
purely imaginary. Therefore the cross term vanishes pointwise:
\begin{equation}\label{eq:kin-split}
 \norm{(\Delta+k^2)u_R}_2^2
 =\norm{\Delta v_R}_2^2+4k^2\norm{\partial_xv_R}_2^2.
\end{equation}
Using the normalized Gaussian moments
\begin{align}\label{eq:moments}
 \norm{\partial_xf_{a,b}}_2^2&=\frac{1}{2a^2},\\
 \norm{\partial_{xx}f_{a,b}}_2^2&=\frac{3}{4a^4},
 &\norm{\partial_{yy}f_{a,b}}_2^2&=\frac{3}{4b^4},\notag\\
 \int_{\R^2}\partial_{xx}f_{a,b}\,\partial_{yy}f_{a,b}
 \,dx\,dy&=\frac{1}{4a^2b^2},\notag
\end{align}
we obtain
\begin{equation}\label{eq:delta-f}
 \norm{\Delta f_{a,b}}_2^2
 =\frac{3}{4a^4}+\frac{1}{2a^2b^2}+\frac{3}{4b^4}.
\end{equation}
The scaling of $v_R$ in dimension two gives
\begin{equation}\label{eq:kinetic}
 \boxed{\begin{aligned}
 \norm{(\Delta+k^2)u_R}_2^2
 &=\frac{2mk^2}{a^2R^2}\\
 &\quad+\frac{m}{R^4}\left(
 \frac{3}{4a^4}+\frac{1}{2a^2b^2}+\frac{3}{4b^4}\right).
 \end{aligned}}
\end{equation}

\medskip
\noindent\textbf{Step 3: The quartic and sextic contributions.}
Direct integration yields
\begin{equation}\label{eq:nonlinear-moments}
 \norm{f_{a,b}}_4^4=\frac{1}{2\pi ab},
 \qquad \norm{f_{a,b}}_6^6=\frac{1}{3\pi^2a^2b^2}.
\end{equation}
Thus the mass-preserving rescaling \eqref{eq:trial} gives
\begin{equation}\label{eq:nonlinear}
 \norm{u_R}_4^4=\frac{m^2}{2\pi abR^2},
 \qquad
 \norm{u_R}_6^6=\frac{m^3}{3\pi^2a^2b^2R^4}.
\end{equation}

\medskip
\noindent\textbf{Step 4: An exact energy polynomial.}
Substitution of \eqref{eq:kinetic} and \eqref{eq:nonlinear}
into \eqref{eq:H} gives the \emph{exact} identity
\begin{equation}\label{eq:polynomial}
 \boxed{H(u_R)=\frac{c_2}{R^2}+\frac{c_4}{R^4},}
\end{equation}
where
\begin{align}\label{eq:c2}
 c_2&=\frac{2mk^2}{a^2}-\frac{\beta m^2}{4\pi ab},\\
 c_4&=m\left(\frac{3}{4a^4}+\frac{1}{2a^2b^2}
 +\frac{3}{4b^4}\right)
 +\frac{\gamma m^3}{9\pi^2a^2b^2}>0.\label{eq:c4}
\end{align}
No higher-order remainder appears in \eqref{eq:polynomial}.

\medskip
\noindent\textbf{Step 5: Choosing the anisotropy.}
From \eqref{eq:c2},
\begin{equation}\label{eq:anisotropy}
 c_2<0
 \quad\Longleftrightarrow\quad
 \frac{a}{b}>\frac{8\pi k^2}{\beta m}.
\end{equation}
For each fixed $m>0$, the right-hand side is finite;
choose any $a,b>0$ satisfying \eqref{eq:anisotropy}.
Then $c_2<0<c_4$.

\medskip
\noindent\textbf{Step 6: Optimizing the spatial scale.}
Setting $z=R^{-2}>0$ in \eqref{eq:polynomial} gives
$H(u_R)=c_2z+c_4z^2$.
Its unique positive minimizer is
\begin{equation}\label{eq:Ropt}
 z_*=-\frac{c_2}{2c_4}>0,
 \qquad R_*^2=-\frac{2c_4}{c_2}>0.
\end{equation}
At this value,
\begin{equation}\label{eq:negative-trial}
 H(u_{R_*})=-\frac{c_2^2}{4c_4}<0.
\end{equation}
Because $\norm{u_{R_*}}_2^2=m$, the definition of the infimum yields
\begin{equation}\label{eq:upper}
 I(m)\le H(u_{R_*})=-\frac{c_2^2}{4c_4}<0.
\end{equation}
This establishes strict negativity for every $m>0$.

\medskip
\noindent\textbf{Step 7: A finite lower bound.}
For every $r\ge0$, the algebraic identity
\begin{equation}\label{eq:completion}
 -\frac{\beta}{2}r^2+\frac{\gamma}{3}r^3
 =\frac{\gamma}{3}r\left(r-\frac{3\beta}{4\gamma}\right)^2
 -\frac{3\beta^2}{16\gamma}r
\end{equation}
shows that
\begin{equation}\label{eq:pointwise}
 -\frac{\beta}{2}|u|^4+\frac{\gamma}{3}|u|^6
 \ge-\frac{3\beta^2}{16\gamma}|u|^2.
\end{equation}
Integration, together with $\norm{(\Delta+k^2)u}_2^2\ge0$, implies
\begin{equation}\label{eq:lower}
 H(u)\ge-\frac{3\beta^2}{16\gamma}\norm{u}_2^2.
\end{equation}
Taking the infimum over the mass-$m$ constraint yields
$I(m)\ge-3\beta^2m/(16\gamma)$.
Combining this estimate with \eqref{eq:upper} proves
\eqref{eq:main}.
\end{proof}

\section{Scope and relation to earlier work}
The plane-wave carrier obeys $(\Delta+k^2)e^{ikx}=0$, eliminating
the constant spectral cost while preserving an admissible
square-integrable Gaussian envelope. Anisotropy makes the coefficient
of $R^{-2}$ negative. The positive defocusing-sextic contribution
appears only in the coefficient of $R^{-4}$ for these trials.
The negative-energy mechanism is already present for the
single-power focusing problem ($\gamma=0$); the new observation here
is its persistence in the combined focusing-cubic/defocusing-quintic
functional with the explicit coefficients above.

The theorem concerns only the value of the constrained infimum. It
neither asserts attainment nor proves strict subadditivity, compactness
of all minimizing sequences, uniqueness of minimizers, or orbital
stability. Such statements require separate arguments.

\begin{thebibliography}{9}
\bibitem{fernandez2022}
Fern\'andez et al., \emph{Journal of Differential Equations}
\textbf{330} (2022),
\href{https://doi.org/10.1016/j.jde.2022.04.037}
{doi:10.1016/j.jde.2022.04.037}.
\end{thebibliography}

\end{document}
